% Anatomía de un diagrama de flujo — la lógica antes que los bloques.
%
% Diagrama propio en TikZ (sin librerías extra: las formas se dibujan a mano).
% Se tiñe con el color de la línea (milcGradeDark = naranja de Robótica). Muestra
% las 4 figuras del diagrama de flujo con un ejemplo real de sensor → umbral →
% acción, para que el estudiante lea la lógica sin ambigüedad antes de programar.
%
% Se incrusta vía \guiaDiagrama desde el bloque `recursos:` del YAML.
\begin{tikzpicture}[scale=1,font=\sffamily,>=stealth,
    every node/.style={font=\scriptsize},
    lbl/.style={font=\scriptsize\itshape,milcNegro!70}]

  % Colores base
  \def\cbox{milcGradeDark}

  % ── (1) Inicio: óvalo ────────────────────────────────────────────────
  \begin{scope}[shift={(0,7.6)}]
    \fill[\cbox] (-1.2,-0.35) rectangle (1.2,0.35);
    \fill[\cbox] (-1.2,0) circle (0.35);
    \fill[\cbox] (1.2,0) circle (0.35);
    \node[white,font=\scriptsize\bfseries] {inicio};
  \end{scope}
  \node[lbl,anchor=west] at (1.9,7.6) {óvalo = empieza / termina};

  % ── (2) Proceso: rectángulo (leer sensor) ───────────────────────────
  \fill[\cbox!85] (-1.5,5.9) rectangle (1.5,6.7);
  \node[white,font=\scriptsize\bfseries,align=center] at (0,6.3) {leer luz\\(sensor)};
  \node[lbl,anchor=west] at (1.9,6.3) {rectángulo = una acción / un paso};

  % ── (3) Decisión: rombo (umbral) ─────────────────────────────────────
  \fill[milcMostaza] (0,5.2) -- (1.7,4.3) -- (0,3.4) -- (-1.7,4.3) -- cycle;
  \draw[milcNegro!55,line width=.4pt] (0,5.2) -- (1.7,4.3) -- (0,3.4) -- (-1.7,4.3) -- cycle;
  \node[milcNegro,font=\scriptsize\bfseries,align=center] at (0,4.3) {¿luz $<$\\umbral?};
  \node[lbl,anchor=west] at (1.9,4.3) {rombo = una pregunta (sí / no)};

  % ── (4) Dos acciones según la respuesta ──────────────────────────────
  \fill[\cbox!85] (-3.3,2.0) rectangle (-0.7,2.9);
  \node[white,font=\scriptsize\bfseries,align=center] at (-2.0,2.45) {mostrar luna\\(oscuro)};

  \fill[\cbox!85] (0.7,2.0) rectangle (3.3,2.9);
  \node[white,font=\scriptsize\bfseries,align=center] at (2.0,2.45) {mostrar sol\\(con luz)};

  % ── (5) Fin ──────────────────────────────────────────────────────────
  \begin{scope}[shift={(0,0.9)}]
    \fill[\cbox] (-1.0,-0.32) rectangle (1.0,0.32);
    \fill[\cbox] (-1.0,0) circle (0.32);
    \fill[\cbox] (1.0,0) circle (0.32);
    \node[white,font=\scriptsize\bfseries] {volver a empezar};
  \end{scope}

  % ── Flechas ──────────────────────────────────────────────────────────
  \draw[->,milcNegro!70,line width=.7pt] (0,7.25) -- (0,6.7);
  \draw[->,milcNegro!70,line width=.7pt] (0,5.9) -- (0,5.2);
  \draw[->,milcNegro!70,line width=.7pt] (-1.7,4.3) -| (-2.0,2.9) node[midway,above,font=\scriptsize\bfseries,milcNegro] {sí};
  \draw[->,milcNegro!70,line width=.7pt] (1.7,4.3) -| (2.0,2.9) node[midway,above,font=\scriptsize\bfseries,milcNegro] {no};
  \draw[->,milcNegro!70,line width=.7pt] (-2.0,2.0) |- (0,1.22);
  \draw[->,milcNegro!70,line width=.7pt] (2.0,2.0) |- (0,1.22);

  % ── Moraleja ─────────────────────────────────────────────────────────
  \node[font=\scriptsize\itshape,milcNegro!75,align=center,text width=11cm]
    at (0.3,0.05)
    {El diagrama de flujo dibuja la lógica con cuatro figuras: \textbf{óvalo}
     (inicio/fin), \textbf{rectángulo} (una acción), \textbf{rombo} (una pregunta
     sí/no) y \textbf{flechas} (el camino). Si se lee claro aquí, los bloques del
     módulo 3 solo lo traducen.};

\end{tikzpicture}
